%%%%%%%%%%%%%%%%%%%%%%%%%%%%%%%%%%%%%%%%%%%%%%%%%%%%%%%%%%%%
% Document settings
\documentclass{ACGSeminar}

%%%%%%%%%%%%%%%%%%%%%%%%%%%%%%%%%%%%%%%%%%%%%%%%%%%%%%%%%%%%
% Own Packages

%%%%%%%%%%%%%%%%%%%%%%%%%%%%%%%%%%%%%%%%%%%%%%%%%%%%%%%%%%%%
% Own Definitions
\newcommand{\comment}[1]{}


%%%%%%%%%%%%%%%%%%%%%%%%%%%%%%%%%%%%%%%%%%%%%%%%%%%%%%%%%%%%
% BibTex
\bibliography{references}

%%%%%%%%%%%%%%%%%%%%%%%%%%%%%
% Hyphenations here
%%%%%%%%%%%%%%%%%%%%%%%%%%%%%
\hyphenation{Sa-tan-arch-aeo-li-deal-co-hell-ish}


%%%%%%%%%%%%%%%%%%%%%%%%%%%%%
% Title, Author, etc.

\begin{document}

\title{Re: Deep G-Buffers for Stable Global Illumination Approximation}

\author{Ferit Tohidi Far}

\maketitle

%%%%%%%%%%%%%%%%%%%%%%%%%%%%%%%%%%%%%%%%%%%%%%%%%%%%%%%%%%%%
% Abstract

\begin{abstract}%
G-Buffers can be used to efficiently render images with a large amount of light sources compared to other local illumination methods. This is possible thanks to a 
process called "deferred rendering". By using Deep G-Buffers we can even approximate global illumination more efficiently than traditional methods like pathtracing. 
\end{abstract}

\keywords{g-buffer, deep g-buffer, pathtracing, global illumination, shading}
\tableofcontents

\newpage

%%%%%%%%%%%%%%%%%%%%%%%%%%%%%%%%%%%%%%%%%%%%%%%%%%%%%%%%%%%%
% Introduction
\label{cha:introduction}
\section{Global Illumination}
	Global Illumination is achieved by taking into account direct light as well as indirect light, which essentially means that it is necessary to compute reflections to a
	certain degree. The most popular method to do this in 3D animated movies is pathtracing (source pixar paper).
	\subsection{Pathtracing}
		
\section{Deferred Rendering}
	\subsection{How deferred rendering handles lighting more efficiently}
	The goal of deferred rendering is to postpone the shading stage. Instead one computes necessary geometry buffers (G-Buffer) and caches them for later use.
	For all practical purposes, G-Buffers have to at least consist of a frame-buffer, a normal-buffer and a z-buffer. Using only these G-Buffers, it is possible
	to compute shading separately.
	\subsection{Deferred Shading}

\section{Geometry Buffer (G-Buffer)}
	\subsection{Frame-Buffer (Image-Buffer)}
	\subsection{Z-Buffer}
	\subsection{Normal-Buffer}
	\subsection{Position-Buffer}
	\subsection{Diffuse-buffer}
	\subsection{Computing local illumination using G-Buffers}
	\subsection{Computing global illumination using G-Buffers}

\section{Deep G-Buffer}
	\subsection{Concept}
	\subsection{How Deep G-Buffers improve performance}

\section{Performance and Output Comparison}
	\subsection{G-Buffers vs Deep G-Buffers vs Pathtracing}

%\section{My Fancy Second Section}
%My second section is even better than the first one\footnote{Proof omitted.} because it contains a %footnote as well as some citations~\cite{Squarepants03,Knuth97}.%
%Everything else is explained in Figure~\ref{fig:dummy}. Always m%ake sure that all figures are refer%enced.
%
%\begin{figure}[htb!]
%  \begin{centering}
%    \includegraphics[width=6cm]{img/thing.jpg}\par
%  \end{centering}
%  \caption{A simple dummy picture to demonstrate how to include graphics into your report.}
%  \label{fig:dummy}
%\end{figure}

%%%%%%%%%%%%%%%%%%%%%%%%%%%%%%%%%%%%%%%%%%%%%%%%%%%%%%%%%%%%
% Bibliography
\label{cha:references}
\printbibliography

\end{document}
